This chapter presents an overview of the statistical analysis used in the $VH$ cross section measurement, including the profile likelihood procedure, the test statistic and corresponding $p$-values, and the neural network binning strategies employed. Results for the indiviudal \wh{} channels, the combined \wh{} channel, the combined Run~3 $VH$ analysis, and the combined Run~2 and Run~3 $VH$ analyses are provided.

Expected results are derived from Asimov datasets in the SRs and observed data in the CRs, while observed results use real data across both signal and control regions. The Asimov dataset is a pseudo-dataset in which the observed event counts in each bin of the signal and control regions are set equal to the expected yields under a specified hypothesis, without any statistical fluctuations. In this analysis, it is constructed under the signal-plus-background hypothesis and is used to estimate the expected significance and uncertainties prior to unblinding. All results account for both statistical and systematic uncertainties across the channels.

For each channel, both expected and observed results include post-fit neural network discriminant distributions, signal strength measurements and significance estimates. Additional outputs will include likelihood scans, impact ranking plots, and correlation matrices of the NPs.

To mitigate large pulls associated with potential mis-modeling of the $WWW$ background as seen in Run~2, the $WWW$ NF is floated in the fit. However, this parameter is only well-constrained in the \zdep{} channel. As a result, post-fit results will only include the $WWW$ NF in the \zdep{} channel and any combinations that include it.

Signal yields and signal strength parameters are obtained via a likelihood fit, described in detail in the following sections. In the combined fit, a simultaneous fit is performed across the neural network discriminant bins of the $VH$ SRs and the total event yields in the CRs. The CRs are used to extract NFs for the dominant prompt backgrounds: $WZ$ for the \wh{} channel and $ZZ$ for the \zh{} channel. Additional fits are performed independently in each SR. The contribution from limited MC statistics is modeled with an additional Poisson term and a NP $\gamma$ for each bin. The $\gamma$ parameter accounts for the statistical uncertainty of the finite MC sample size by allowing predicted yield in the bin to fluctuate within its Poisson uncertainty.